\section{Data Hasil Pengamatan dan Ekstraksi}

\subsection*{Dataset Eksperimen IAEA-NDS EXFOR Inti \texorpdfstring{$^{12}\text{C}(n,\text{tot})$}{12C(n,tot)}}
Dataset eksperimen yang digunakan bersumber dari repositori internasional \textit{International Atomic Energy Agency Nuclear Data Services} (IAEA-NDS) EXFOR entri 10224002 berdasarkan pengukuran transmisi resolusi tinggi oleh \textcite{cierjacks1968high}. Rentang energi kinetik neutron yang dianalisis dibatasi secara fisis pada $1.8000\text{ MeV} \le E \le 2.4000\text{ MeV}$ yang memuat struktur resonansi tunggal gelombang-D.

Dataset numerik terdiri atas $N = 59$ titik observasi presisi ganda (\texttt{float64}) yang diekspor ke dalam lembar data \texttt{exfor\_c12\_resonance.csv} dan dapat diakses pada repositori GitHub \url{https://github.com/xysilverberg/Neutron-Resonance-Breit-Wigner}. Cuplikan 20 titik data eksperimen representatif di sekitar puncak resonansi disajikan pada Tabel~\ref{tab:data_exfor}.

\begin{table}[H]
    \centering
    \footnotesize
    \setlength{\tabcolsep}{6pt}
    \caption{Cuplikan data eksperimen penampang lintang $^{12}\text{C}(n,\text{tot})$ dari repositori EXFOR.}
    \label{tab:data_exfor}
    \begin{tabular}{c S[table-format=1.4] S[table-format=1.4] S[table-format=1.4]}
        \toprule
        {No.} & {Energi Neutron $E_i$ (\si{\mega\electronvolt})} & {Penampang Lintang $\sigma_i$ (barn)} & {Ketidakpastian $\Delta\sigma_i$ (barn)} \\
        \midrule
        1  & 1.8000 & 4.6520 & 0.0750 \\
        5  & 1.8400 & 4.6730 & 0.0780 \\
        10 & 1.8900 & 4.6880 & 0.0760 \\
        15 & 1.9400 & 4.7520 & 0.0800 \\
        18 & 1.9700 & 4.8320 & 0.0830 \\
        20 & 1.9900 & 4.9560 & 0.0860 \\
        22 & 2.0100 & 5.2050 & 0.0900 \\
        24 & 2.0300 & 5.7480 & 0.0950 \\
        25 & 2.0400 & 6.2300 & 0.1010 \\
        26 & 2.0500 & 6.8450 & 0.1080 \\
        27 & 2.0600 & 7.3820 & 0.1140 \\
        28 & 2.0700 & 7.5950 & 0.1180 \\
        29 & 2.0750 & 7.5810 & 0.1170 \\
        30 & 2.0780 & 7.5720 & 0.1160 \\
        31 & 2.0800 & 7.5350 & 0.1150 \\
        33 & 2.0900 & 7.1850 & 0.1110 \\
        35 & 2.1100 & 6.0450 & 0.0980 \\
        40 & 2.1600 & 4.9120 & 0.0840 \\
        48 & 2.2400 & 4.7120 & 0.0760 \\
        59 & 2.4000 & 4.6620 & 0.0710 \\
        \bottomrule
    \end{tabular}
\end{table}

\subsection*{Ekstraksi Otomatis Tebakan Parameter Awal (\texorpdfstring{$\mathbf{p}_0$}{p0})}
Menggunakan modul heuristik fisis pada Tahap 3, diperoleh nilai estimasi awal parameter sebelum loop Newton-Raphson dijalankan:
\begin{enumerate}
    \item \textbf{Energi Resonansi Awal}: $E_r^{(0)} = E_{\arg\max(\sigma)} = 2.0700\text{ MeV}$.
    \item \textbf{Kontinum Baseline Latar Belakang}: $\sigma_{bg}^{(0)} = \min(\sigma_1, \sigma_N) = 4.6520\text{ barn}$.
    \item \textbf{Tinggi Amplitudo Puncak Murni}: $\sigma_0^{(0)} = \sigma_{\max} - \sigma_{bg}^{(0)} = 7.5950 - 4.6520 = 2.9430\text{ barn}$.
    \item \textbf{Lebar Resonansi FWHM}: Nilai ambang batas setengah tinggi adalah:
    \begin{equation}
        \sigma_{\text{half}} = \sigma_{bg}^{(0)} + 0.5\,\sigma_0^{(0)} = 4.6520 + 0.5(2.9430) = 6.1235\text{ barn}
    \end{equation}
    Rentang energi di atas nilai tersebut adalah $E \in [2.0400, 2.1000]\text{ MeV}$, sehingga diperoleh lebar awal $\Gamma^{(0)} = 2.1000 - 2.0400 = 0.0600\text{ MeV}$.
\end{enumerate}

Dengan vektor tebakan awal $\mathbf{p}_0 = [2.0700, 0.0600, 2.9430, 4.6520]^T$, nilai fungsi objektif awal terhitung sebesar $\chi^2(\mathbf{p}_0) = 90.8117$.
